\kbvchapter{ಪಾಠ ೨೨}{ದಿನಬಳಕೆಯ ಗಾದೆಗಳು}

\noindent\hspace*{0.6em}\kbvbullet\hspace{0.55em}ಅಕ್ಕಿ ಮೇಲೆ ಆಸೆ ನೆಂಟರ ಮೇಲೆ ಇಷ್ಟ.\par

\noindent\hspace*{0.6em}\kbvbullet\hspace{0.55em}ಅತ್ತೆಗೊಂದು ಕಾಲ ಸೊಸೆಗೊಂದು ಕಾಲ.\par

\noindent\hspace*{0.6em}\kbvbullet\hspace{0.55em}ಅತಿ ಆಸೆ ಗತಿ ಕೆಡಿಸಿತು.\par

\noindent\hspace*{0.6em}\kbvbullet\hspace{0.55em}ಅಲ್ಪರ ಸಂಗ ಅಭಿಮಾನ ಭಂಗ.\par

\noindent\hspace*{0.6em}\kbvbullet\hspace{0.55em}ಅಗ್ಗದ ಶೆಟ್ಟಿ ಮುಗ್ಗಿದ ಜೋಳ.\par

\noindent\hspace*{0.6em}\kbvbullet\hspace{0.55em}ಅಜ್ಜಿ ನೂತದ್ದೆಲ್ಲಾ ಅಜ್ಜನ ಉಡಿದಾರಕ್ಕೆ.\par

\noindent\hspace*{0.6em}\kbvbullet\hspace{0.55em}ಅನಿಷ್ಟಕ್ಕೆ ಶನೇಶ್ವರ ಗುರಿ,\par

\noindent\hspace*{0.6em}\kbvbullet\hspace{0.55em}ಅಹಂಕಾರಕ್ಕೆ ಉದಾಸೀನವೇ ಮದ್ದು.\par

\noindent\hspace*{0.6em}\kbvbullet\hspace{0.55em}ಅಟ್ಟ ಹತ್ತದವನು ಬೆಟ್ಟ ಹತ್ತಿಯಾನೆ?\par

\noindent\hspace*{0.6em}\kbvbullet\hspace{0.55em}ಅಲ್ಪ ವಿದ್ಯಾ ಮಹಾಗರ್ವಿ.\par

\noindent\hspace*{0.6em}\kbvbullet\hspace{0.55em}ಅಲ್ಪನಿಗೆ ಐಶ್ವರ್ಯ ಬಂದರೆ ನಡುರಾತ್ರೀಲಿ ಕೊಡೆ ಹಿಡಿಸಿಕೊಂಡ.\par

\noindent\hspace*{0.6em}\kbvbullet\hspace{0.55em}ಅಂಗೈ ಹುಣ್ಣಿಗೆ ಕನ್ನಡಿಯೇಕೆ?\par

\noindent\hspace*{0.6em}\kbvbullet\hspace{0.55em}ಆರು ದೋಸೆ ಕೊಟ್ಟರೆ ಅತ್ತೆ ಕಡೆ ಮೂರು ದೋಸೆ ಕೊಟ್ಟರೆ ಸೊಸೆ ಕಡೆ.\par

\noindent\hspace*{0.6em}\kbvbullet\hspace{0.55em}ಆಕಳು ಕಪ್ಪಾದರೆ ಹಾಲು ಕಪ್ಪೇ?\par

\noindent\hspace*{0.6em}\kbvbullet\hspace{0.55em}ಆಳಾಗಬಲ್ಲವನು ಅರಸಾಗಬಲ್ಲ.\par

\noindent\hspace*{0.6em}\kbvbullet\hspace{0.55em}ಆಪತ್ತಿಗೆ ಆದವನೇ ನೆಂಟ.\par

\noindent\hspace*{0.6em}\kbvbullet\hspace{0.55em}ಇಬ್ಬರ ಜಗಳದಲ್ಲಿ ಮೂರನೆಯವನಿಗೆ ಲಾಭ.\par

\noindent\hspace*{0.6em}\kbvbullet\hspace{0.55em}ಇರುಳು ಕಂಡ ಬಾವೀಲಿ ಹಗಲು ಬಿದ್ದ.\par

\noindent\hspace*{0.6em}\kbvbullet\hspace{0.55em}ಇಲಿಗೆ ಪ್ರಾಣ ಸಂಕಟ, ಬೆಕ್ಕಿಗೆ ಚೆಲ್ಲಾಟ.\par

\noindent\hspace*{0.6em}\kbvbullet\hspace{0.55em}ಉರಿಯೋ ಹುಣ್ಣಿನ ಮೇಲೆ ಬರೆ ಹಾಕಿದ ಹಾಗೆ.\par

\noindent\hspace*{0.6em}\kbvbullet\hspace{0.55em}ಉಗುರಿನಲ್ಲಿ ಹೋಗೋದಕ್ಕೆ ಕೊಡಲಿ ಪೆಟ್ಟು.\par

\noindent\hspace*{0.6em}\kbvbullet\hspace{0.55em}ಉಪ್ಪು ತಿಂದವನು ನೀರು ಕುಡೀತಾನೆ.\par

\noindent\hspace*{0.6em}\kbvbullet\hspace{0.55em}ಉಪ್ಪುಂಡ ಮನೆಗೆ ಎರಡು ಬಗೆಯಬಾರದು.\par

\noindent\hspace*{0.6em}\kbvbullet\hspace{0.55em}ಊರಿದ್ದ ಕಡೆ ಹಗರಿ,\par

\noindent\hspace*{0.6em}\kbvbullet\hspace{0.55em}ಊಟ ಬಲ್ಲವನಿಗೆ ರೋಗವಿಲ್ಲ, ಮಾತು ಬಲ್ಲವನಿಗೆ ಜಗಳವಿಲ್ಲ.\par
